\chapter{The C Model: Memory and Control Flow}

\textbf{Readings:} CP:AMA 5.2--5.3, 6, 7.1--7.3, 9.2--9.3, 18.2.

\section{Models of Computation}

In CS 135 we modelled Racket's behaviour with \emph{substitution}. Substitution cannot handle
mutation, so in CS 136 we model C with two components:
\begin{itemize}
  \item \textbf{memory}: what data is stored where, and
  \item \textbf{control flow}: which statement is being executed, and what runs next.
\end{itemize}

\section{Memory}

\subsection{Bits, bytes and addresses}

A \emph{bit} of storage has two possible states: 0 or 1. A \emph{byte} is 8 bits, so a byte has
$2^8 = 256$ possible states. Memory is a huge array of bytes; the position of a byte is its
\emph{address}. To access data we go to its address. Addresses are usually written in
hexadecimal, e.g.\ \code{0x5000}.

\subsection{\code{sizeof}}

When C defines a variable it reserves enough consecutive bytes for its type. The \code{sizeof}
operator (it looks like a function but is an operator evaluated at compile time) gives the number
of bytes a type or expression uses. Its result has type \code{size\_t}, printed with \code{\%zd}.
\begin{lstlisting}[style=c]
printf("%zd\n", sizeof(int));    // 4 (in Seashell)
printf("%zd\n", sizeof(char));   // 1 (always)
int n = 0;
printf("%zd\n", sizeof n);       // 4
\end{lstlisting}

\subsection{Integer limits and overflow}

An \code{int} in Seashell is 4 bytes $=$ 32 bits, so there are only $2^{32}$ possible values.
Integers are \emph{signed}: the range is $-2^{31}$ to $2^{31}-1$. The limits are provided by
\code{<limits.h>}:
\begin{lstlisting}[style=c]
#include <limits.h>
INT_MAX   //  2147483647  =  2^31 - 1
INT_MIN   // -2147483648  = -2^31
\end{lstlisting}
(You do not need to memorise these values.)

\begin{important}[Overflow]
If an arithmetic result is outside the range of its type, \emph{overflow} occurs. For signed
\code{int}, overflow is \textbf{undefined behaviour}: anything may happen (in practice the
value usually ``wraps around'', e.g.\ \code{INT\_MAX + 1} becomes \code{INT\_MIN}). It is your
job to make sure overflow cannot happen.
\end{important}

\begin{aside}
\code{unsigned int} uses the same 32 bits to store $0$ to $2^{32}-1$, and overflow of unsigned
values is well defined (arithmetic modulo $2^{32}$). \code{long long} is at least 64 bits.
\end{aside}

\subsection{The \code{char} type and ASCII}

A \code{char} is 1 byte and is used to store a single character. It is really a small integer:
characters are encoded using ASCII (American Standard Code for Information Interchange), which
assigns codes 0--127, e.g.\ \code{'A'} is 65, \code{'a'} is 97, \code{'0'} is 48, \code{' '} is
32 and \code{'\textbackslash n'} is 10.
\begin{lstlisting}[style=c]
char c = 'A';       // the same value as
char d = 65;        // ... this
printf("%c %d\n", c, c);         // A 65
char lower = c + ('a' - 'A');    // 'a'
int digit = '7' - '0';           // 7
\end{lstlisting}
The placeholder for \code{char} is \code{\%c}. Single quotes give a character; double quotes give
a string (Unit~08). Many assignment problems are solved by doing arithmetic on characters as
above.

\subsection{Structures in memory}

A structure occupies (at least) the sum of the sizes of its fields; the compiler may add
\emph{padding} between fields for alignment, so \code{sizeof(struct s)} can be larger than the
sum. Always use \code{sizeof} rather than computing sizes by hand.

\section{Sections of Memory}

The memory available to a C program is divided into five \emph{sections} (also called
segments):

\begin{center}
\begin{tikzpicture}[sec/.style={draw,minimum width=4.2cm,minimum height=0.75cm}]
  \node[sec] (code) {code};
  \node[sec,below=0pt of code] (ro) {read-only data};
  \node[sec,below=0pt of ro] (glob) {global data};
  \node[sec,below=0pt of glob] (heap) {heap $\downarrow$};
  \node[sec,below=0pt of heap,minimum height=1.2cm,fill=gray!10] (gap) {(free)};
  \node[sec,below=0pt of gap] (stack) {stack $\uparrow$};
  \node[left=0.3cm of code] {low addresses};
  \node[left=0.3cm of stack] {high addresses};
\end{tikzpicture}
\end{center}

The operating system and hardware enforce the boundaries: trying to access memory outside the
sections your program owns (or writing to a read-only section) causes a
\textbf{segmentation fault}, and the program crashes.

\subsection{The code section}

\emph{Source code} is the human-readable ASCII text we write. The compiler translates it into
\emph{machine code}, which the processor can execute; this translation is called
\emph{compiling}. When the program runs, the machine code is placed in the code section.
(CS 241 shows how this is done.)

\subsection{Read-only and global data}

Global (top-level) identifiers are stored in one of these two sections:
\begin{itemize}
  \item global \emph{constants} go in the \textbf{read-only data} section;
  \item global \emph{variables} go in the \textbf{global data} section.
\end{itemize}
These are created and initialized \emph{before} \code{main} runs, and exist until the program
ends. This is why global initializers must be constant expressions.

The heap is covered in Unit~10, and the stack below.

\section{Control Flow}

\subsection{Return addresses and the call stack}

When a function is called, the program ``jumps'' to the function's code. When it finishes,
execution continues where the call was made. So for each call we must remember the
\emph{return address}: the location in the caller to go back to.

\begin{remark}
The return address is not the same as the \code{return} statement. \code{return} produces a
value for the caller; the return address says \emph{where} in the caller to resume. Even
\code{main} has a return address (back into the operating system).
\end{remark}

Calls are nested: \code{main} calls \code{f}, which calls \code{g}, \dots, and they return in the
reverse order. So the history of calls is a \emph{stack} (LIFO), called the \emph{call stack}. It
is the same idea as the stack ADT: a call pushes, a return pops.

\subsection{Stack frames}

Each function call pushes a \emph{stack frame} (activation record) containing:
\begin{itemize}
  \item the argument values (parameters),
  \item the local variables,
  \item the return address.
\end{itemize}
Each call gets a new frame, whether it is a call to a different function or a recursive call to
the same function.

\begin{important}[Pass by value]
C makes a \textbf{copy} of each argument value and stores the copy in the new stack frame. This
is called \emph{pass by value}. Changing a parameter inside a function does not change the
caller's variable. (Unit~06 shows how to emulate \emph{pass by reference} with pointers.)
\end{important}

\begin{example}
\begin{lstlisting}[style=c]
int fact(int n) {
  if (n <= 1) return 1;
  return n * fact(n - 1);   // fact:7
}

int main(void) {
  int r = fact(3);          // main:11
  printf("%d\n", r);
}
\end{lstlisting}
When \code{fact(1)} is executing, the call stack looks like this (top of the stack drawn at the
top):
\begin{center}
\begin{tikzpicture}[fr/.style={draw,minimum width=5cm,align=left,text width=4.8cm,inner sep=4pt}]
  \node[fr] (f1) {\textbf{fact}\\ n: 1\\ return address: fact:7};
  \node[fr,below=0pt of f1] (f2) {\textbf{fact}\\ n: 2\\ return address: fact:7};
  \node[fr,below=0pt of f2] (f3) {\textbf{fact}\\ n: 3\\ return address: main:11};
  \node[fr,below=0pt of f3] (m) {\textbf{main}\\ r: ???\\ return address: OS};
\end{tikzpicture}
\end{center}
\end{example}

\subsection{The stack section}

The call stack is stored in the stack section of memory. Each frame uses memory, so very deep
recursion (or very large local variables) can exhaust it: a \textbf{stack overflow}, which usually
crashes the program with a segmentation fault. The stack section in Seashell is relatively small
(a few megabytes), so recursion depth of around a million is already a problem.

\section{The Model}

At any moment, the state of a running C program is:
\begin{enumerate}
  \item the \emph{current program location} (which statement is being executed), and
  \item the current contents of memory.
\end{enumerate}

\paragraph{Calling a function.}
\begin{enumerate}
  \item A stack frame is created and pushed; argument values are copied into it.
  \item The current program location is stored in the frame as the return address.
  \item The program location jumps to the start of the function.
\end{enumerate}
Local variables are added to the frame as their definitions are reached.

\paragraph{Returning.}
\begin{enumerate}
  \item The return value (if any) is noted.
  \item The program location is changed back to the frame's return address.
  \item The stack frame is popped (its memory is no longer valid).
\end{enumerate}

\section{Control Flow Statements}

\subsection{\code{if}}
\begin{lstlisting}[style=c]
if (condition) {
  // block
}

if (condition) {
  // block
} else {
  // block
}

if (c1) {
  // block
} else if (c2) {
  // block
} else {
  // block
}
\end{lstlisting}
Braces may be omitted for a single statement, but that is error-prone (``dangling else''); CS 136
style always uses braces.

\subsection{Looping}

With mutation we can repeat a block until some state changes, which replaces most uses of
recursion.

\paragraph{\code{while}.}
\begin{lstlisting}[style=c]
// sum_first(n) produces 1 + ... + n
int sum_first(int n) {
  int sum = 0;
  int i = 1;
  while (i <= n) {
    sum += i;
    ++i;
  }
  return sum;
}
\end{lstlisting}
The condition is checked before each iteration; if it is false initially the body never runs.
A common bug is forgetting to update the loop variable, giving an infinite loop.

\paragraph{\code{do ... while}.} The body runs at least once, and the condition is checked
after each iteration:
\begin{lstlisting}[style=c]
do {
  // body
} while (condition);
\end{lstlisting}

\paragraph{\code{for}.} The most common loop combines initialization, condition and update:
\begin{lstlisting}[style=c]
for (setup; condition; update) {
  body;
}
// is equivalent to
setup;
while (condition) {
  body;
  update;
}
\end{lstlisting}
\begin{lstlisting}[style=c]
int sum = 0;
for (int i = 1; i <= n; ++i) {   // i is local to the loop (C99)
  sum += i;
}
\end{lstlisting}

\paragraph{Nested loops.}
\begin{lstlisting}[style=c]
// print_triangle(n) prints a right triangle of stars with n rows
// effects: produces output
void print_triangle(int n) {
  for (int row = 1; row <= n; ++row) {
    for (int col = 0; col < row; ++col) {
      printf("*");
    }
    printf("\n");
  }
}
\end{lstlisting}

\paragraph{\code{break} and \code{continue}.} \code{break} exits the innermost loop immediately;
\code{continue} skips the rest of the current iteration and goes to the next one (in a
\code{for} loop, the update still runs).
\begin{lstlisting}[style=c]
// is_prime(n) determines if n is prime
bool is_prime(int n) {
  if (n < 2) return false;
  for (int d = 2; d * d <= n; ++d) {
    if (n % d == 0) {
      return false;          // return also exits the loop
    }
  }
  return true;
}
\end{lstlisting}

\subsection{Recursion vs.\ iteration}

Accumulative recursion translates directly into a loop: the accumulators become variables that
are mutated each iteration. A loop uses a single stack frame, so it cannot overflow the stack.

\section{Summary}
\begin{itemize}
  \item Memory is an array of bytes with addresses; \code{sizeof} gives sizes in bytes.
  \item \code{int} has a limited range (\code{INT\_MIN} to \code{INT\_MAX}); signed overflow is
    undefined behaviour. \code{char} is a 1-byte integer holding an ASCII code.
  \item Five sections: code, read-only data, global data, heap, stack. Bad access causes a
    segmentation fault.
  \item Each call pushes a stack frame (arguments, locals, return address); arguments are passed
    by value. Too many frames cause a stack overflow.
  \item \code{if}, \code{while}, \code{do}--\code{while}, \code{for}, \code{break},
    \code{continue}.
\end{itemize}
